\section{Tests et validation}
\label{sec:tests}

\subsection{Stratégie de tests}

La stratégie repose sur trois niveaux complémentaires~:

\begin{description}
    \item[Tests unitaires] couvrent les méthodes métier pures, en
    particulier le calcul du ROI sur \texttt{Campaign}, le service
    \texttt{predict\_potential()} et la construction du contexte CRM
    (\texttt{build\_crm\_context}).
    \item[Tests d'intégration] couvrent les vues HTTP avec le
    \texttt{Client} de Django (création de produit, déplacement Kanban,
    soumission du formulaire de prédiction). L'appel DeepSeek est
    \emph{mocké} via \texttt{responses} pour rester déterministe.
    \item[Tests manuels] focus sur l'i18n (chaque page chargée en EN,
    FR, AR avec vérification visuelle du RTL) et sur la fluidité du
    chatbot en streaming.
\end{description}

\subsection{Exemple de test unitaire ROI}

\begin{lstlisting}[language=Python,caption={tests/test\_roi.py}]
import pytest
from decimal import Decimal
from apps.pipeline.models import Product, Campaign
from apps.finance.models import Expense, Revenue

@pytest.mark.django_db
def test_roi_basic():
    p = Product.objects.create(name='Test')
    c = Campaign.objects.create(product=p, name='C1',
        start_date='2026-01-01', budget=1000)
    Expense.objects.create(campaign=c, amount=Decimal('500'),
        category='ads', date='2026-01-02')
    Revenue.objects.create(campaign=c, amount=Decimal('800'),
        date='2026-01-10')
    assert c.roi() == 60  # (800-500)/500 * 100
\end{lstlisting}

\subsection{Validation du modèle ML}

Le script \texttt{ml/train.py} affiche un \texttt{classification\_report}
de scikit-learn sur le jeu de test (20~\% du dataset synthétique). Les
métriques cibles sont~:

\begin{itemize}
    \item Accuracy globale~$\geq 0{,}80$.
    \item F1-score par classe (\emph{Low} / \emph{Medium} / \emph{High})~%
    $\geq 0{,}75$.
    \item Matrice de confusion exportée en PNG dans
    \texttt{docs/report/figures/confusion\_matrix.png}.
\end{itemize}

\subsection{Validation de l'i18n}

Pour chaque langue, on vérifie qu'il n'existe \emph{aucune} entrée fuzzy
ni vide dans les fichiers \texttt{.po}~:

\begin{lstlisting}[language=bash]
django-admin makemessages -a
msgattrib --untranslated locale/fr/LC_MESSAGES/django.po
msgattrib --untranslated locale/ar/LC_MESSAGES/django.po
\end{lstlisting}

Aucune sortie attendue.

\subsection{Validation du chatbot}

\begin{enumerate}
    \item Test de connectivité~: appel direct à \texttt{ask\_deepseek()}
    avec un message connu, vérification du code HTTP 200 et de la présence
    d'un champ \texttt{content} non vide.
    \item Test de contextualisation~: création de campagnes fictives,
    appel du builder, vérification que les noms apparaissent bien dans le
    \texttt{system\_prompt}.
    \item Test multilingue~: même question posée en EN, FR, AR~; on vérifie
    que la réponse est dans la langue demandée.
\end{enumerate}
